% !TEX root = ./main.tex
\chapter{Altri Servizi Backend}
\label{app:backend-services}

\section{Panoramica}
Oltre al servizio RAG, il backend di SOULFRAME si regge su tre componenti più specializzati: Whisper per la trascrizione, Coqui XTTS v2 per la sintesi vocale e Avatar Asset Server per import, cache e serving dei modelli. La loro separazione conta perché ciascuno porta con sé vincoli diversi di latenza, I/O e stato persistente.

\par\begingroup
\small
\setlength{\LTleft}{0pt}
\setlength{\LTright}{0pt}
\captionsetup{type=table}
\captionof{table}{Tre servizi backend di supporto al flusso conversazionale.}
\label{tab:appendix-backend-overview}
\begin{longtable}{|>{\raggedright\arraybackslash}p{0.22\textwidth}|>{\raggedright\arraybackslash}p{0.14\textwidth}|>{\raggedright\arraybackslash}p{\dimexpr0.64\textwidth-6\tabcolsep-4\arrayrulewidth\relax}|}
\hline
\textbf{Servizio} & \textbf{Porta} & \textbf{Ruolo principale} \\
\hline
\endfirsthead
\hline
\textbf{Servizio} & \textbf{Porta} & \textbf{Ruolo principale} \\
\hline
\endhead
Whisper & \texttt{8001} & Riceve audio in upload e restituisce la trascrizione. \\
\hline
Coqui XTTS v2 & \texttt{8004} & Gestisce voce per-avatar, wait phrases, sintesi completa e streaming. \\
\hline
Avatar Asset Server & \texttt{8003} & Importa modelli \texttt{.glb}, li deduplica, li conserva in cache e li serve al client. \\
\hline
\end{longtable}
\par\endgroup

Questi tre servizi fanno meno cose del RAG, ma non sono affatto accessori. Whisper apre il canale vocale dell'utente, Coqui chiude il turno con una voce coerente con l'avatar e l'Avatar Asset Server mantiene stabile il catalogo dei modelli tra import, cache e caricamento lato client. La loro separazione non nasce da un gusto astratto per i micro-servizi: nasce dal fatto che hanno tempi, failure mode e costi molto diversi, quindi conviene tenerli osservabili e sostituibili senza trascinare con sé tutto il resto.

\section{Pipeline audio e risposta vocale}
\subsubsection*{Whisper: trascrizione e punto di ingresso del turno.}
Il servizio \texttt{whisper\_server.py} è volutamente compatto. Carica il modello definito in \texttt{WHISPER\_MODEL}, espone \texttt{/health} e accetta upload multipart sulla route \texttt{/transcribe}. Il file viene scritto in una posizione temporanea, trascritto e poi rimosso.

Il contratto che il client usa davvero è minimale:

\begin{lstlisting}[basicstyle=\ttfamily\small,breaklines=true,caption={Campi usati da \texttt{/transcribe}.}]
POST /transcribe
- file: UploadFile
- language: Form (default "it")
\end{lstlisting}

\begin{lstlisting}[language=json,caption={Risposta tipica di \texttt{/transcribe}.}]
{
  "text": "Domani ho un colloquio alle dieci"
}
\end{lstlisting}

\begin{lstlisting}[language=json,caption={Risposta tipica di \texttt{/health} nel servizio Whisper.}]
{
  "ok": true,
  "model": "small"
}
\end{lstlisting}

Il servizio non fa molto altro, e questa è una scelta positiva. Whisper qui si comporta come un ingresso pulito del turno: riceve l'audio, lo trascrive e restituisce testo. I controlli sulla durata minima dell'input, sullo scarto di trascrizioni spurie e sulla gestione dell'attesa rimangono nel client, che è il punto in cui quegli errori diventano davvero percepibili. In questo modo lo STT non viene caricato di logiche che appartengono più alla UX del turno che alla trascrizione in sé.

\paragraph{Coqui XTTS v2: voce per-avatar, streaming e wait phrases.}
Il servizio \texttt{coqui\_tts\_server.py} è più articolato. Oltre alle route di sintesi, gestisce il reference vocale per-avatar, la generazione delle clip di attesa e il caricamento anticipato del modello.

\paragraph{Profili vocali.}
Il reference vocale viene salvato tramite \texttt{/set\_avatar\_voice}. Il file caricato viene letto, convertito in mono, normalizzato e rifinito prima di diventare il \texttt{reference.wav} dell'avatar.

\begin{lstlisting}[basicstyle=\ttfamily\small,breaklines=true,caption={Campi essenziali usati dal client nel setup voce.}]
POST /set_avatar_voice
- avatar_id
- speaker_wav
\end{lstlisting}

Questa fase conta più di quanto sembri, perché definisce la materia prima di tutto il sottosistema vocale. Una reference sporca o troppo debole non peggiora soltanto la qualità timbrica: rende instabile anche l'effetto di continuità tra avatar, wait phrases e risposta principale. Per questo il servizio pulisce il file prima di salvarlo, invece di limitarsi ad archiviarlo così com'è.

\paragraph{Streaming.}
La route \texttt{/tts\_stream} restituisce il reply come \texttt{audio/wav} progressivo. Il backend non espone al client i segmenti interni della sintesi come unità semantiche separate: prepara il reply, lo segmenta solo per gestire meglio la generazione XTTS e poi ricompone il risultato in un unico stream WAV. Il contratto reale, però, non è mono-canale. Accanto allo stream audio esiste un secondo canale logico, basato su \texttt{request\_id} e \texttt{/tts\_timing}, che fornisce snapshot temporali incrementali per il reveal del testo lato client.

Dal lato client contano soprattutto due campi: \texttt{request\_id}, usato per correlare audio e timing, e \texttt{client\_platform}, che permette al backend di distinguere il percorso WebGL da quello nativo. Se il client non è WebGL e il \texttt{request\_id} è presente, il servizio apre una sessione temporale incrementale e la aggiorna man mano che i segmenti del reply vengono sintetizzati e allineati. Se invece il client è WebGL, la stessa sessione non viene attivata: il reply continua a uscire come audio progressivo, ma il canale temporale fine viene soppresso.

\begin{figure}[t]
\centering
\resizebox{0.82\linewidth}{!}{%
\begin{tikzpicture}[
  >=stealth,
  every node/.style={font=\footnotesize, align=center},
  box/.style={draw, rounded corners, minimum height=1.1cm, text width=2.8cm, inner sep=5pt},
  lbl/.style={font=\scriptsize, fill=white, inner sep=1.5pt}
]
  \node[box] (reply) at (0,0) {Reply testuale};
  \node[box, text width=3.2cm] (seg) at (4.4,0) {Segmentazione\\progressiva XTTS};
  \node[box, text width=3cm] (audio) at (9.2,1.55) {Stream audio\\\texttt{/tts\_stream}};
  \node[box, text width=3cm] (timing) at (9.2,-1.55) {Timing snapshot\\\texttt{/tts\_timing}};
  \node[box, text width=3.3cm] (native) at (14.4,1.55) {Client nativo\\reveal temporizzato};
  \node[box, text width=3.3cm] (webgl) at (14.4,-1.55) {Client WebGL\\fallback statico};
  \draw[->] (reply) -- node[lbl, above] {\texttt{request\_id}} (seg);
  \draw[->] (seg) -- (audio);
  \draw[->] (seg) -- node[lbl, above right] {snapshot cumulativi} (timing);
  \draw[->] (audio) -- (native);
  \draw[->] (timing) -- (native);
  \draw[->] (audio) -- (webgl);
  \draw[->] (timing) -- node[lbl, below] {disattivato} (webgl);
\end{tikzpicture}
}
\captionsetup{width=0.92\linewidth}
\caption{Schema concettuale del Word Flow: il reply viene segmentato per la sintesi, lo stream audio parte subito, mentre i timing incrementali restano attivi solo nel percorso nativo. In WebGL il testo non viene fatto avanzare parola per parola, ma mostrato direttamente in forma completa.}
\label{fig:wordflow-concept}
\end{figure}

Lo stream espone alcuni header operativi utili al client:

\par\begingroup
\small
\setlength{\LTleft}{0pt}
\setlength{\LTright}{0pt}
\captionsetup{type=table}
\captionof{table}{Header operativi restituiti da \texttt{/tts\_stream}.}
\label{tab:appendix-tts-stream-headers}
\begin{longtable}{|>{\raggedright\arraybackslash}p{0.24\textwidth}|>{\raggedright\arraybackslash}p{\dimexpr0.76\textwidth-4\tabcolsep-3\arrayrulewidth\relax}|}
\hline
\textbf{Header} & \textbf{Uso nel flusso} \\
\hline
\endfirsthead
\hline
\textbf{Header} & \textbf{Uso nel flusso} \\
\hline
\endhead
\texttt{X-Avatar-Id} & Rende esplicito quale profilo vocale sia stato usato. \\
\hline
\texttt{X-Coqui-Lang} & Conferma la lingua effettivamente adottata nello stream. \\
\hline
\texttt{X-Audio-Rate} & Espone il sample rate del flusso audio. \\
\hline
\texttt{X-Audio-Format} & Indica che il contenuto è PCM16 dentro un contenitore WAV. \\
\hline
\texttt{X-Speaker-Source} & Distingue se la voce proviene da reference salvata, upload temporaneo o speaker di default. \\
\hline
\texttt{X-TTS-Request-Id} & Identificatore della risposta TTS, usato dal client per interrogare i tempi incrementali su \texttt{/tts\_timing}. \\
\hline
\end{longtable}
\par\endgroup

La logica reale del reply tiene insieme tre passaggi distinti. Prima il testo viene pulito e suddiviso in segmenti compatibili con XTTS, usando punteggiatura forte e un limite di caratteri per evitare truncation sui blocchi troppo lunghi. Poi ogni segmento viene sintetizzato in WAV e, quando disponibile, allineato parola-per-parola con \emph{forced alignment}. Infine il backend invia subito l'audio in streaming progressivo e aggiorna in parallelo una sessione di timing incrementale associata al \texttt{request\_id}. Il frontend riceve quindi un solo flusso audio da \texttt{/tts\_stream} e interroga periodicamente \texttt{/tts\_timing} per ottenere \texttt{words}, \texttt{word\_end\_ms} e \texttt{segment\_end\_ms} cumulativi. Il punto importante è che questi array sono cumulativi sull'intero reply, non riservati al singolo segmento corrente: così il client può far avanzare il testo senza ricostruire localmente una timeline alternativa.

Nel contratto corrente i parametri opzionali più rilevanti lato client sono \texttt{reply\_segment\_max\_chars}, usato per limitare la lunghezza massima dei segmenti del reply prima della sintesi, e \texttt{request\_id}, usato per correlare stream audio e snapshot temporali. I riferimenti più vecchi a \texttt{split\_sentences}, \texttt{max\_chunk\_chars} e ai vecchi header temporali non descrivono più il comportamento effettivo del reply in \texttt{MainMode}.

Quando l'allineamento non è disponibile, il backend continua comunque a restituire l'audio del reply. In quel caso la sessione \texttt{/tts\_timing} rimane vuota o incompleta, con eventuale campo \texttt{error}, e il client degrada a una resa testuale statica invece di ricostruire stime locali poco affidabili. In WebGL questo fallback non è occasionale ma sistematico: il \emph{word flow} temporizzato del reply è disabilitato e il testo viene mostrato direttamente in forma completa. La ragione è operativa più che estetica: nel runtime Unity WebGL le chiamate ripetute necessarie a tenere insieme audio progressivo, timing incrementale e sincronizzazione fine con il layer visivo hanno mostrato un costo eccessivo rispetto al beneficio percepito. Una soluzione più robusta richiederebbe una revisione più profonda del sottosistema di sincronizzazione.

\paragraph{Wait phrases.}
Le clip di attesa vengono preparate quando il client salva il riferimento vocale dell'avatar con \path{/set_avatar_voice}. In quella fase il backend usa lo stesso \texttt{reference.wav} appena validato per sintetizzare un piccolo insieme di file brevi (\texttt{wait\_<key>.wav}) e salvarli nella directory dell'avatar. Durante la conversazione, invece, \path{/wait_phrase} non sintetizza nulla: restituisce semplicemente una delle clip già presenti. Questa scelta tiene separate due responsabilità: il setup voce costruisce il materiale necessario, mentre il turno conversazionale si limita a richiamarlo in modo rapido e prevedibile. È utile tenerla distinta anche dalla risposta principale: le wait phrases non usano \texttt{request\_id}, non aprono una sessione \texttt{/tts\_timing} e non partecipano al \emph{word flow} del reply.

\begin{lstlisting}[language=json,caption={Risposta tipica di \texttt{/set\_avatar\_voice}.}]
{
  "ok": true,
  "avatar_id": "LOCAL_model1",
  "bytes": 184320,
  "wait_phrases_ready": true
}
\end{lstlisting}

Qui si nota una differenza importante rispetto alla risposta principale. Le wait phrases non devono essere semanticamente ricche: devono solo riempire l'attesa senza stonare rispetto al timbro dell'avatar. Per questo stanno nello stesso servizio della TTS vera e propria, ma con una funzione diversa: non costruiscono il contenuto del turno, ne proteggono il ritmo.

\section{Gestione asset avatar}
\subsubsection*{Avatar Asset Server: import, cache e catalogo.}
Il servizio \texttt{avatar\_asset\_server.py} costruisce il catalogo avatar lato backend. Non si limita a fare da file server: riceve URL di export, scarica il modello, deduplica l'asset tramite hash dell'URL sorgente, aggiorna il file di metadati e restituisce un \texttt{cached\_glb\_url} coerente con l'ambiente di esecuzione.

\begin{lstlisting}[language=json,caption={Payload tipico della route \texttt{/avatars/import}.}]
{
  "avatar_id": "019c8296-1a77-75b1-9d4b-c36ff20d1a4c",
  "url": "https://api.avaturn.me/avatars/exports/.../model",
  "gender": "male",
  "bodyId": "body_1",
  "urlType": "httpURL"
}
\end{lstlisting}

L'output tipico è il seguente:

\begin{lstlisting}[language=json,caption={Campi principali restituiti da \texttt{/avatars/import}.}]
{
  "ok": true,
  "avatar_id": "019c8296-1a77-75b1-9d4b-c36ff20d1a4c",
  "cached_glb_url": "/api/avatar/avatars/019c8296-1a77-75b1-9d4b-c36ff20d1a4c/model.glb",
  "bytes": 11983396,
  "dedup": false
}
\end{lstlisting}

La route \texttt{/avatars/list} unifica poi:
\begin{itemize}
\item modelli locali di fallback;
\item avatar importati e mantenuti nello store;
\item URL già normalizzati per il client.
\end{itemize}

\begin{lstlisting}[language=json,caption={Forma sintetica della risposta di \texttt{/avatars/list}.}]
{
  "avatars": [
    {
      "avatar_id": "LOCAL_model1",
      "source": "local"
    },
    {
      "avatar_id": "019c8296-1a77-75b1-9d4b-c36ff20d1a4c",
      "source": "avaturn",
      "cached_glb_url": "/api/avatar/avatars/019c8296-1a77-75b1-9d4b-c36ff20d1a4c/model.glb"
    }
  ]
}
\end{lstlisting}

Nel codice c'è anche un aspetto poco vistoso ma utile: durante \texttt{/avatars/list}, gli elementi il cui file \texttt{.glb} non è più risolvibile vengono scartati e il file di metadati viene riallineato. Il catalogo esegue quindi una correzione automatica e limitata dei metadati\footnote{Per correzione automatica e limitata dei metadati si intende qui il riallineamento del catalogo quando esso rileva voci ormai non risolvibili.}, così il client evita di ricevere voci già rotte in partenza. Questo dettaglio conta perché la libreria avatar non è un semplice menu: è il punto in cui persistenza server, cache e stato del frontend devono restare coerenti anche dopo import ripetuti, cancellazioni o riavvii.

La separazione tra questi tre servizi resta utile perché hanno costi, failure mode e tempi molto diversi. Whisper è essenzialmente un ingresso vocale con upload temporaneo; Coqui è il punto più sensibile alla latenza di risposta; l'Avatar Asset Server è soprattutto I/O e persistenza. Tenerli distinti rende più semplice il debugging e impedisce che un problema di un sottosistema renda opaco l'intero backend. Ognuno resta così abbastanza piccolo da essere leggibile, ma anche abbastanza autonomo da poter essere osservato e sostituito senza riscrivere il resto della pipeline.
